\section{总结与展望}

本文围绕智能结构优化设计场景下的等几何分析研究，对统一几何-分析描述、结构优化经典数学模型以及密度法、水平集法、MMC 和多补丁壳体优化等主要技术路线进行了综述。综合全文，主要认识如下。

（1）等几何分析的核心价值并不局限于高阶离散本身，而在于以统一样条模型贯通 CAD 建模、结构分析和优化更新，使拓扑优化、形状优化与厚度优化能够共享同一几何基础。

（2）从方法与应用进展看，等几何结构优化已由规则连续体问题扩展到多补丁壳体、壳型胞元、超材料以及与增材制造协同的复杂轻量化结构；同时，数据驱动设计、智能生成和几何回写也在重塑结构优化的实现路径。

（3）当前研究仍存在复杂几何参数化、局部细化、多尺度建模、可制造性约束以及数据驱动与物理模型耦合不足等问题。后续研究应进一步建立开放基准、统一数据格式和可复现工作流，使不同算法能够在相同几何、载荷、约束和制造条件下进行比较。

从工程应用角度看，等几何分析的价值最终要落实到设计流程中，即优化结果不仅应具有较优性能，还应服务于概念设计、方案细化、制造约束校核和几何回写。随着开源工作流、自动微分工具和智能生成模型逐步成熟，等几何分析有望进一步发展为可嵌入工程设计平台的通用能力。
